% ---------- Слайд 11: \pause ----------
\begin{frame}[fragile,label=pause]{Поэтапный вывод: \cmd{pause}}
  \begin{columns}[T,onlytextwidth]
    \column{0.52\textwidth}
      \textbf{Исходный код:}
      \vspace{0.2em}
      \inputminted[fontsize=\scriptsize, frame=single, breaklines=true]{latex}{examples/pause.tex}

    \column{0.44\textwidth}
      \textbf{Результат (листайте слайды):}
      \vspace{0.2em}
      \begin{block}{}
        Этап 1: постановка задачи; \\ \pause
Этап 2: метод решения; \\ \pause
Этап 3: результаты.
      \end{block}
  \end{columns}
  \note{Типичная ошибка --- показать весь слайд сразу. \cmd{pause} синхронизирует восприятие.}
\end{frame}

% ---------- Слайд 12: Спецификаторы <...> ----------
\begin{frame}[fragile,label=overlays]{Управление видимостью: \texttt{<...>}}
  \begin{columns}[T,onlytextwidth]
    % --- ЛЕВАЯ КОЛОНКА: КОД ---
    \column{0.52\textwidth}
      \textbf{Исходный код:}
      \vspace{0.2em}
      \inputminted[fontsize=\scriptsize, frame=single, breaklines=true]{latex}{examples/overlays.tex}

    % --- ПРАВАЯ КОЛОНКА: РЕЗУЛЬТАТ ---
    \column{0.44\textwidth}
      \textbf{Результат (листайте слайды):}
      \vspace{0.2em}
      \begin{block}{}
        \begin{itemize}
  \item<3-> Появляется на шаге 3;
  \item<2-> Появляется на шаге 2;
  \item<1-> Появляется на шаге 1.
\end{itemize}
      \end{block}
  \end{columns}
  \note{Спецификаторы <...> обеспечивают точечное управление
  видимостью. Например, \cmd{alert<2>\{текст\}} выделяет текст
  только на втором шаге.}
\end{frame}